\begin{abstract}
Mutually unbiased bases describe quantum measurements for which certainty in one basis gives uniform outcomes in every other. Their maximum number is known when the dimension is a power of a prime number, but the smallest unresolved case, dimension six, has resisted both constructions and exclusion proofs for decades. Here we present a computer-assisted proof that this maximum is three. We use our complete classification of complex Hadamard matrices, which allowed a reconstruction based on four phases,   with exact Fourier identities and certified local exclusions. The maxima of a scalar function allow us to extend local exclusions throughout one part of the reconstruction. An exact identity forces the remaining candidates into the Fourier class, which is excluded. 
We also determine the phase regions occupied by the known triples and conjecture that a two-parameter construction and an isolated Fourier triple exhaust them. Dimension six thus attains the tensor product bound, even though constructions in dimension twelve exceed it. This fixes the number of exactly complementary projective measurements in a six-level system, while the classification of the triples remains open. 
\end{abstract}
\maketitle

\section{Introduction}
How many quantum measurements can be mutually complementary? Mutually unbiased bases give a precise version of this question. For $d\ge2$, an orthonormal basis of a $d$-dimensional quantum system is a set of $d$ mutually orthogonal, normalized states. These states label the outcomes of a projective measurement. Let $\mathcal B=\{\ket{b_j}\}_{j=1}^d$ and $\mathcal C=\{\ket{c_k}\}_{k=1}^d$ be two bases of $\mathbb{C}^d$. They are \textit{mutually unbiased} when
\begin{equation}
 \abs{\braket{b_j}{c_k}}^2=\frac1d
 \quad\text{for all }j,k.
 \label{eq:mub}
\end{equation}
Preparing any state of one basis then gives equal probabilities for all outcomes in the other. A set of mutually unbiased bases (MUBs) satisfies this condition for every pair. The open problem is to determine the largest possible set for a given dimension.

An order-$d$ complex Hadamard matrix is a matrix $H\in\C^{d\times d}$ with $|H_{jk}|=1$ and $H^\dagger H=dI_d$. If $A$ and $B$ contain the vectors of two orthonormal bases as columns, Eq.~\eqref{eq:mub} says exactly that $\sqrt d\,A^\dagger B$ is Hadamard. Here $H^\dagger$ is the conjugate transpose and $I_d$ is the identity matrix. 
%Hadamard matrices are unnormalized throughout, whereas matrices whose columns are basis vectors are unitary.

MUBs grew from Schwinger's work on unitary operators and Ivanovi\'c's study of quantum state determination~\cite{Schwinger1960,Ivanovic1981}. Later, Wootters and Fields showed how MUBs give optimal state reconstruction~\cite{WoottersFields1989}. They also enter uncertainty relations~\cite{MaassenUffink1988}, quantum key distribution~\cite{Cerf2002}, and tests of entanglement~\cite{Spengler2012}. Write $N(d)$ for their maximum number. One always has $N(d)\le d+1$~\mcw{[CITE]}. A set of $d+1$ bases is called complete, and such sets exist whenever $d$ is an integer power of a prime. Outside prime-power dimensions, no upper bound has even been found~\cite{McNultyWeigert2026}. 

For $d=\prod_{i=1}^s p_i^{a_i}$, with distinct primes $p_i$ and positive integers $a_i$, tensoring complete MUB sets in the prime-power factors gives the \emph{tensor product bound}~\cite{Klappenecker2004} $N(d)\ge 1+\min_i p_i^{a_i}$, which is a lower bound where each product basis uses one basis from every factor, so the factor with the fewest bases limits this construction.

\begin{figure*}[t]
 \centering
 \includegraphics[width=0.8\textwidth]{Figure_1.jpg}
 \caption{\textbf{Two reductions behind the no four MUB proof.} \textbf{a, Matched-completion reduction.} The large torus $\T^2_{12}$ retains the first two corner angles. Blue circles (A), triangles (B), and stars (C) mark three groups of six maximum seeds; circles and triangles coincide in this projection. The smaller tori $\T^2_{34}$ resolve the other two angles at the indicated fixed $(\theta_1/\alpha,\theta_2/\alpha)$, with $\alpha=\arccos(1/4)$. The large surface shows $M=\max_{\theta_3,\theta_4}(R/R_{\max})$; the small surfaces show slices of $R/R_{\max}$, with nonpositive values gray. All eighteen maxima yield one matched-completion class.
 \textbf{b, All-critical reduction.} The separate Fourier torus parameterizes whole Fourier-family matrices. Its scale $Q$ is their normalized sum of $R$ over all three-by-three corners. The twelve orange zeros are Fourier-equivalent; the all-critical classification reduces to this class. Both reductions lead to excluded classes. The shared colorbar refers to the three quantities labeled in the figure, and faded symbols lie on hidden surfaces. These plots illustrate the exact results, rather than supply the exclusion certificate; the slices and scales are specified in SM Sec.~II.}
 \label{fig:proof}
\end{figure*}

Dimension six is the smallest case outside the prime powers, and it has remained open for decades since Zauner conjectured in his dissertation that $N(6)=3$~\cite{Zauner1999}. The MUB existence problem is one of the \emph{Five Open Problems in Quantum Information Theory} discussed by Horodecki, Rudnicki, and \.{Z}yczkowski~\cite{HorodeckiRudnickiZyczkowski2022}, \mcw{ADD MORE DESCRIPTIONS} 
Three MUBs can easily be obtained by combining three MUBs of a two-level system with three of a three-level system gives three bases in $\mathbb C^2\otimes\mathbb C^3$. Such bases have been implemented with photons, using polarization and orbital angular momentum~\cite{DAmbrosio2013}. 

Rigorous exclusions concern specified pairs or triples: Grassl's triples~\cite{Grassl2004}, triples made entirely of product states~\cite{McNultyWeigert2012Triples}, and pairs containing the standard basis and a member of the two-parameter Fourier family~\cite{Jaming2009}. The latter result combines analytic error bounds with exhaustive computation. Tao's exceptional pair cannot even extend to a third basis~\cite{BrierleyWeigert2009}. These results do not cover every pair, so a proof of $N(6)=3$ requires excluding arbitrary quartets. 


Bengtsson and collaborators mapped triples through order-six Hadamard matrices~\cite{BengtssonEtAl2007}. Numerical searches over bases and smaller constellations~\cite{ButterleyHall2007,BrierleyWeigert2008}, and Bell-inequality optimizations~\cite{PratColomer2022}, support $N(6)=3$. Tests modulo primes~\cite{McConnell2021} and polynomial optimization~\cite{Mortimer2023} give further restrictions. Recently, we showed that a phased unitary construction can give five MUBs in dimension twelve~\cite{wuttig2026mutuallyunbiasedbases}, exceeding the tensor product bound of four even in small dimensions where numerical searches \mcw{[CITE]} did not violate the tensor product bound. The tensor product bound is therefore not always sharp, making its status in dimension six a separate question.

Building on Sz\"oll\H{o}si's four-parameter construction~\cite{Szollosi2012FourParameter}, our complete classification reconstructs every order-six Hadamard, up to row and column phases and permutations, from a suitable dephased $3\times3$ corner~\cite{wuttig2026completeclassificationcomplexhadamard}. Dephasing leaves four corner phases. At a suitable corner, the compatible right and lower blocks form finite lists whose pairings recover the matrix, including the Karlsson family and Tao's isolated matrix. Except for Tao's class and one explicit Karlsson matrix, these lists follow from one quadratic and one cubic in each direction.

Here, we prove $N(6)=3$ by first excluding all conjugate-label matched completions through the maxima of a scalar function and certified local exclusions. An exact nonnegative identity forces any remaining extendable matrix into the already excluded Fourier class. We then locate the known triples in phase coordinates and state their proposed classification as Conjecture~\ref{conj:triples}. \mcw{REMOVE? Appendix~\ref{app:constellations} gives exact partial constellations; the Outlook discusses the positive minimum error for approximate quartets.}

\noindent\emph{Related approaches.}
During the development of this work, we became aware of OpenAI's reported computer-assisted proof of $N(6)=3$, based on a nine-phase cover~\cite{OAIMUB2026}. We have not independently verified the full proof, as the computational certificates needed to reproduce the reported exclusion were unavailable in the release we consulted. 
Our proof is based on a reconstruction of all complex Hadamard matrices~\cite{wuttig2026completeclassificationcomplexhadamard}, which reduces the continuous description from nine covering coordinates to four corner phases. A global maximum argument then extends certified local exclusions to every matched completion, while our complementary-invariant identity and all-critical classification reduce the remaining candidates to the excluded Fourier class. These structural reductions remove the need for a global nine-phase cover, confining interval computations to local exclusions and boundary transfers. Our work uses the single-matrix Fourier-vanishing theorem and exact moment-reduction machinery in the complementary-invariant step~\cite{OAIFourier2026}, and adapts the phase-ball geometry and local interval-exclusion framework~\cite{OAIMUB2026}. \mcw{This has to be written more understandable and easier to read} Our local certificates and verification records are specified in the SUpplementary Material.

\section{No four mutually unbiased bases}
\label{sec:nofour}
Choose one hypothetical basis as the coordinate basis $U_0=I_6$. Call the pair $I_6,H/\sqrt6$ \emph{extendable} if two further bases can be added to form four MUBs. The other bases would be $U_j=H_j/\sqrt6$, with each $H_j$ Hadamard. Their relative transitions must also be Hadamard, since
\[
 U_j^\dagger U_k=\frac{H_j^\dagger H_k}{6},\qquad
 \sqrt6\,U_j^\dagger U_k=\frac{H_j^\dagger H_k}{\sqrt6}.
\]
Being unbiased to $I_6$ alone does not impose these pairwise conditions. We prove that no extendable pair exists, and hence $N(6)=3$. First we exclude all completions with matched conjugate side labels; then we show that any remaining extendable matrix must be Fourier-equivalent and is also excluded.

\paragraph*{Four-phase reconstruction.}
Row and column phases and permutations define Hadamard equivalence. Extendability is also preserved by conjugation and transposition, which we treat as separate symmetries. Write $\T=\{z\in\C:\abs{z}=1\}$. Dephasing a $3\times3$ corner fixes its first row and column to one, leaving four phases:
\begin{equation}
 E(\theta)=\begin{pmatrix}1&1&1\\1&e^{i\theta_1}&e^{i\theta_2}\\1&e^{i\theta_3}&e^{i\theta_4}\end{pmatrix}.
 \label{eq:seed}
\end{equation}
This \emph{seed} lies in $\T^4=\T^2_{12}\times\T^2_{34}$, where the factors retain the first and last two angles. It specifies a corner, with chart and completion choices still required to obtain a whole matrix.

Write $H=\left(\begin{smallmatrix}E&B\\C&D\end{smallmatrix}\right)$. Block orthogonality gives
\[
 \begin{aligned}
 BB^\dagger&=6I_3-EE^\dagger,& C^\dagger C&=6I_3-E^\dagger E,\\
 CE^\dagger&+DB^\dagger=0.&&
 \end{aligned}
\]
The first two matrices record the side overlaps. With $G=6I_3-EE^\dagger$, set $s=G_{21}$, $t=G_{31}$, $r=G_{32}$, $X=|s|^2$, $Y=|t|^2$, $Z=|r|^2$, $p=X+Y+Z$ and $J=2\operatorname{Re}(sr\bar t)$. Their common invariants give the real scalar
\begin{equation}
 \begin{aligned}
 R(E)&=4g^2\eta-\chi^2,\\
 g&=27-3p+J,\qquad \eta=9-p+J,\\
 \chi&=XY+XZ+YZ-XYZ\\
 &\quad-(p-9)(p-18)+(p-15)J.
 \end{aligned}
 \label{eq:scalar}
\end{equation}
This is the polynomial of SM Eq.~(S.1.1). For a \emph{strict} seed, $R>0$, both Gram matrices are positive definite, and each side admits exactly the two distinct conjugate labels $z_\pm$. A prescribed label fixes an existing side up to phases and permutations (SM Proposition~S.1). Inverting $B^\dagger$ in the last equation forces
\begin{equation}
 \begin{aligned}
 D&=-CE^\dagger(B^\dagger)^{-1}\\
  &=-CE^\dagger(6I_3-EE^\dagger)^{-1}B.
 \end{aligned}
 \label{eq:forced}
\end{equation}
Orthogonality fixes $D$; its entrywise unit modulus is a separate result. SM Proposition~S.2 proves this flatness for the \emph{matched} label pairs $(z_+,z_-)$ and $(z_-,z_+)$ of $B,C^T$, as opposed to equal-label pairs. Both matched completions exist at every strict seed. For the first matching, the equivalence class $f_+(\theta)=[H_+(\theta)]$ varies continuously, including across auxiliary root cuts. Conjugating the seed and completion exchanges the matchings.

Every Hadamard class other than Tao's and the all-critical Fourier class occurs among these matched strict completions. Indeed, as shown below, absence of a matched strict corner implies Eq.~\eqref{eq:complement}, which forces a non-Tao matrix to be Fourier-equivalent. These two exceptional classes are treated separately.

\paragraph*{Excluding the matched completions.}
Work in the quotient of order-six Hadamards by row and column phases and permutations, with the topology inherited from the phase entries. The certificate of SM Proposition~S.4 covers actual Hadamards in $11\,228$ closed phase neighborhoods. Outward-rounded interval bounds enclose all common unbiased vectors and exclude the compatibility pattern required for two additional bases, uniformly including each neighborhood's boundary. Let $W$ be the union of these quotient images and their conjugates, transposes, and conjugate transposes. Every member is nonextendable. The images are compact, so their finite union is closed. Strict boundary-transfer inequalities place every boundary Hadamard in a covered interior, making $W$ open. Records $0$ and $877$ contain interior witnesses for a strict matched completion and the Fourier matrix, respectively.

Set $\Omega=\{\theta:R(\theta)>0\}$ and $A=f_+^{-1}(W)$. Continuity makes $A$ open and closed relative to $\Omega$, and the strict witness makes it nonempty. Suppose $\Omega\setminus A$ were also nonempty. Maximize $R$ on the closure of each set in the compact torus $\T^4$. Both maxima are positive, so their maximizing seeds remain in $\Omega$. Relative closedness puts each seed back in its original set; openness makes it a local maximum of $R$. SM Proposition~S.5 classifies \emph{all} positive local maxima as the eighteen-seed orbit of
\begin{equation}
 E_* =\begin{pmatrix}1&1&1\\1&x&x^2\\1&x^{-1}&x^{-2}\end{pmatrix},
 \qquad x=\frac{1+i\sqrt{15}}4.
 \label{eq:maximum}
\end{equation}
Here $x=e^{i\alpha}$, with $\alpha=\arccos(1/4)$ as in Fig.~\ref{fig:proof}. Row and column permutations followed by dephasing generate the orbit and preserve the side labels; side uniqueness identifies all eighteen $f_+$ outputs with one class. That class cannot be both inside and outside $W$. Thus $A=\Omega$, without assuming that $\Omega$ is connected. Conjugation invariance of $W$ excludes the other matching as well.

\paragraph*{All other candidates reduce to Fourier.}
Suppose an extendable Hadamard $H$ remains. It lies outside $W$, so its sides cannot have matched labels at any strict corner. To compare complementary blocks independently of their row and column phases, define, for any phase block $X$,
\begin{equation}
 \Psi(X)=\frac{\perm(X)\det(X)^2}{\prod_{a,b}X_{ab}}.
 \label{eq:psi}
\end{equation}
The permanent has the determinant's terms without alternating signs. Each row or column phase contributes the same cubic factor to numerator and denominator, so $\Psi$ is phase invariant.

For actual sides $B,C^T$, the labels obey $z=-\Psi/g$, where $g=|\det B|^2=|\det C|^2$. Equal labels therefore give $\Psi(B)=\Psi(C)$. At $R=0$ the labels coalesce when $g>0$; when $g=0$, both invariants vanish. Actual corners cannot have $R<0$, by the side identity in SM Eq.~(S.1.2). Thus these algebraic identities cover nonstrict and singular corners, including the Karlsson exception to the generic recipe. Exchanging column blocks gives, by SM Proposition~S.6,
\begin{equation}
 \Psi(H[I,J])=\Psi(H[I^c,J^c]),\qquad |I|=|J|=3.
 \label{eq:complement}
\end{equation}
This holds for every choice of three rows $I$ and three columns $J$; superscript $c$ denotes the complementary indices. Tao's class is separately excluded from triples~\cite{BrierleyWeigert2009}.

Measure the failure of complementary invariants to be complex conjugates by
\begin{equation}
 \mathcal D(H)=\frac1{400}\sum_{|I|=|J|=3}
 \left|\Psi(H[I,J])-\overline{\Psi(H[I^c,J^c])}\right|^2,
 \label{eq:defect}
\end{equation}
where $400=\binom63^2$. For an ordered three-by-three submatrix $X$, form
\[
 \frac{X_{12}X_{13}X_{21}X_{23}X_{31}X_{32}}
 {(X_{11}X_{22}X_{33})^2}.
\]
Define $S(H)$ as the average of this ratio over the $120$ ordered row triples and $120$ ordered column triples.

For every non-Tao Hadamard, the imported Fourier theorem~\cite{OAIFourier2026} makes the mean over columns of the product of any three rows and the conjugates of the other three vanish, and likewise for $H^T$ (SM Eq.~(S.2.2)). Under these vanishings, Hadamard orthogonality and Eq.~\eqref{eq:complement}, our exact rational moment reduction (SM Proposition~S.7) gives
\begin{equation}
 \mathcal D(H)+\bigl(30\operatorname{Re}S(H)-3\bigr)^2=0.
 \label{eq:sos}
\end{equation}
Both terms are nonnegative, so
\[
 \mathcal D(H)=0\ \Longrightarrow\
 \Psi(H[I,J])=\overline{\Psi(H[I^c,J^c])}\quad\forall I,J.
\]
Equality in Eq.~\eqref{eq:complement} and conjugate equality together force every $\Psi$ to be real. For each actual seed $E$ and its adjacent side $B$, SM Proposition~S.1 gives $R(E)=4(\operatorname{Im}\Psi(B))^2$. Since $B$ is itself a three-by-three submatrix, $R(E)=0$ at every corner: $H$ is \emph{all-critical}. This side identity applies to actual sides, not arbitrary points of $\T^4$.

SM Proposition~S.3 states that every all-critical order-six Hadamard is Fourier-equivalent, including singular and exceptional cases. The Fourier class is in $W$, contradicting extendability. Thus no arbitrary quartet exists and $N(6)\le3$; the explicit triple in Eq.~\eqref{eq:product} gives $N(6)\ge3$.

The almost unbiased quartet of Raynal, L\"u, and Englert~\cite{Raynal2011} attains average squared basis distance approximately $0.9983$, where one means exact mutual unbiasedness. Its coordinate basis is unbiased to the other three, but the remaining pairwise errors are nonzero. Our proof does not establish that this quartet is the closest possible one.

\begin{figure*}[t]
 \centering
 \includegraphics[width=\textwidth]{Figure_2.jpg}
\caption{\textbf{Where the known triples appear.}
\textbf{a,} Triangle edges identify the three pairwise transitions, ordered as the columns in \textbf{c--e}.
\textbf{b,} The square shows corner-angle coordinates of Sz\"oll\H{o}si's two-circulant family $X^{(2)}$, flattening the lower torus in \textbf{c} with opposite edges identified.
\textbf{c--e,} Each transition's four corner angles are split between a lower $(\theta_1,\theta_2)$ torus and an upper $(\theta_3,\theta_4)$ torus. Green marks pairs reached by the continuous construction; dotted gray marks pairs not reached in that chart. Blue traces the Zauner subfamily; orange stars $p_\pm$ mark the isolated Fourier triple and its conjugate. Insets magnify the upper $p_+$ projections: the continuous construction and Zauner subfamily project to a line in \textbf{d} and a point in \textbf{e}. Coincident projections need not represent equivalent triples. Equation~\eqref{eq:regions} specifies the images; Conjecture~\ref{conj:triples} asserts exhaustiveness of the known constructions.}
 \label{fig:triples}
\end{figure*}

\section{The geometry of three bases}
\label{sec:triples}
Three MUBs do exist in dimension six. The simplest example takes tensor products of triples in dimensions two and three:
\begin{equation}
 \begin{aligned}
 U_0&=I_6,\qquad U_1=F_2\otimes F_3,\\
 U_2&=(D_2F_2)\otimes(D_3F_3),\\
 D_2&=\diag(1,i),\qquad D_3=\diag(1,\omega,\omega),
 \end{aligned}
 \label{eq:product}
\end{equation}
Here $F_n$ is the normalized Fourier matrix, with entries $(F_n)_{ab}=n^{-1/2}e^{2\pi iab/n}$ for $a,b=0,\ldots,n-1$, and $\omega=e^{2\pi i/3}$. The diagonal matrices $D_2,D_3$ supply the phases in the second product basis. The squared overlaps in the two factors are $1/2$ and $1/3$, so their products are $1/6$. This is the isolated Fourier triple. There is also a continuous two-parameter construction due to Zauner and Sz\"oll\H{o}si~\cite{Zauner1999,Szollosi2010Hypocycloids,Jaming2009}. We ask whether these account for every triple, and where they appear in the four-phase coordinates used in the proof.

\paragraph*{How the continuous triples are made.}
Keep $U_0=I_6$ and choose the second basis
\begin{equation}
 U_1=\frac1{\sqrt2}
 \begin{pmatrix}F_3&DF_3\\F_3&-DF_3\end{pmatrix},
 \label{eq:Qmain}
\end{equation}
where $D=\diag(1,e^{ia},e^{ib})$ contains two free phases. For every $a,b$, there is a vector $v=(r,s)^T/\sqrt6$ with unit-modulus entries in $r,s$ that is unbiased to both $U_0,U_1$~\cite{IdelWolf2015}. Set $Z=\diag(1,\omega,\omega^2,1,\omega,\omega^2)$ and $J_D(r,s)^T=(D\bar s,-D\bar r)^T$. The six vectors $Z^kv,Z^kJ_Dv$, $k=0,1,2$, remain common unbiased vectors and are mutually orthogonal; they form $U_2$. SM Eq.~(S.3.1) and its following calculation verify this special construction. Different starting vectors can give different triples, and every admissible choice is included.

The triangle in Fig.~\ref{fig:triples} organizes the three transitions $H_{jk}=\sqrt6\,U_j^\dagger U_k$. Their product telescopes:
\begin{equation}
 \begin{aligned}
 H_{01}H_{12}H_{20}
 &=(\sqrt6)^3U_0^\dagger U_1U_1^\dagger U_2U_2^\dagger U_0\\
 &=6\sqrt6\,I_6.
 \end{aligned}
 \label{eq:cycle}
\end{equation}
Two triples are equivalent under a common unitary, phases and permutations of basis vectors, and permutation of the bases. Within-basis changes give $H_{jk}\mapsto M_j^\dagger H_{jk}M_k$, where the $M_j$ are monomial unitaries: the shared factors cancel around the triangle. Independent Hadamard equivalences need not preserve this product constraint.

\refstepcounter{conjecture}\label{conj:triples}
\paragraph*{Conjecture \theconjecture{} (triple classification~\cite{MatolcsiEtAl2026MUBTriplets}).}
Every MUB triple in $\C^6$ is equivalent either to a triple from the continuous construction above, with arbitrary $a,b$ and every admissible choice of the starting vector, or to the isolated triple in Eq.~\eqref{eq:product}.

For the continuous construction, the transition families can be ordered as
\[
 H_{12}\in X^{(2)},\qquad
 H_{01}\in(F_6^{(2)})^T,\qquad H_{20}\in F_6^{(2)},
\]
where membership is up to Hadamard equivalence. Here $X^{(2)}$ is Sz\"oll\H{o}si's family built from two $3\times3$ circulant blocks, in which each row is a cyclic shift of the preceding row. The two-parameter Fourier family $F_6^{(2)}$ consists, up to equivalence, of the matrices $\sqrt6\,U_1^T$ from Eq.~\eqref{eq:Qmain} as $a,b$ vary. For the isolated triple, all three transitions are equivalent to the unnormalized Fourier matrix $\sqrt6F_6$.

\paragraph*{What the tori show.}
Each column of Fig.~\ref{fig:triples} follows one transition as the continuous construction varies. Its dephased corner has four angles: the lower torus retains $(\theta_1,\theta_2)$ and the upper torus $(\theta_3,\theta_4)$.

To obtain the displayed corners, we use $H_{12}$, $H_{01}P$, and $P^TH_{20}$, where $P$ reorders columns as $(0,3,1,4,2,5)$. These are choices of chart for the individual matrices. Equation~\eqref{eq:cycle} applies to the transitions before these chart reorderings. Let $G$ denote the allowed angle pairs for the two circulant blocks, described below. The four seed angles reached by the construction are
\begin{equation}
 \begin{aligned}
 H_{12}:\quad &(u,v,u-v,u),\quad (u,v)\in G,\\
 H_{01}P:\quad &(a,2\pi/3,b,-2\pi/3),\\
 P^TH_{20}:\quad &(c,d,-2\pi/3,2\pi/3).
 \end{aligned}
 \label{eq:regions}
\end{equation}
Within the separate transition images, $(a,b)$ and $(c,d)$ each range over $[0,2\pi)^2$; they are not four freely combinable parameters of one triple. All angles are modulo $2\pi$. The first column maps $(u,v)\in G$ to $(u-v,u)$ above; the middle column has two circle projections; the right has a full lower torus and the single upper point $(-2\pi/3,2\pi/3)$.

The boundary of $G$ expresses orthogonality of the two circulant blocks. The first supplies three unit phases; the second must supply three more, of product one, whose sum cancels the first sum. Such phases exist exactly in $G$, including its boundary (SM Eq.~(S.3.4)). These are transition coordinates, whereas $(a,b)$ in Eq.~\eqref{eq:Qmain} specify the second basis; their regions need not have the same shape.

Green records the continuous construction's exact image in each chosen chart, not a classification of all triples. A corner omits completion data, each torus discards two angles, and the three transition images remain coupled by Eq.~\eqref{eq:cycle}. Thus an orange Fourier star in green indicates coincident projected coordinates, not equivalence of full triples or membership in the continuous family. Conjugation can move a representative outside the green chart while preserving its Hadamard class. The stars $p_\pm$ mark the Fourier representative and its conjugate; blue follows the Zauner subfamily.

Every triple in Conjecture~\ref{conj:triples} has a Fourier-family transition. The known obstruction for such a pair~\cite{Jaming2009} already prevents a fourth basis. Proving the conjecture would therefore give another route to $N(6)=3$. Our upper bound establishes the number of bases without assuming that the known triples exhaust them.

\section{Outlook}
The equality $N(6)=3$ makes every MUB triple in dimension six unextendible by a fourth basis and rules out a complete set of seven. The same bound applies to collections of complementary commuting observables. By their correspondence with MUBs~\cite{Boykin2007}, the traceless $6\times6$ complex matrices contain at most three pairwise orthogonal Cartan subalgebras that are closed under Hermitian conjugation. These are maximal commuting subalgebras whose matrices can be diagonalized together. Orthogonality uses the Killing form, which here is proportional to the trace of the matrix product. In particular, a decomposition into seven such subalgebras is impossible. Closure under Hermitian conjugation is essential to this consequence, so the unrestricted Lie-algebra problem remains open.

Three MUB measurements supply $3(6-1)=15$ independent probabilities, whereas a general six-level state has $6^2-1=35$ real parameters. Tomography therefore needs additional measurements, or an informationally complete generalized measurement~\cite{WoottersFields1989,Scott2006}. Cryptographic and entanglement tests can use three exact MUBs or larger approximate sets, accounting for the measured overlaps.

A symmetric informationally complete positive operator-valued measurement (SIC-POVM) in dimension six uses $36$ unit vectors with $\abs{\braket{\psi_j}{\psi_k}}^2=1/7$ for $j\ne k$, and effects $\dyad{\psi_j}/6$~\cite{Renes2004}. Grassl constructed such a SIC exactly~\cite{Grassl2004}. Its vectors form a complex projective $2$-design, reproducing uniformly distributed pure states' second moments, and its probabilities determine an arbitrary state~\cite{Renes2004,Scott2006}. Thus SIC tomography remains possible without seven MUBs; Zauner's SIC conjecture in general dimension remains a separate existence question~\cite{Zauner1999}.

Proving Conjecture~\ref{conj:triples} would classify the maximal sets. For approximate quartets, define the total overlap error
\[
 C=\sum_{0\le j<k\le3}\sum_{a,b=1}^6
 \left(|(U_j^\dagger U_k)_{ab}|^2-\frac16\right)^2.
\]
Compactness of $U(6)^4$ and continuity give $C_{\min}>0$, since zero would be an exact quartet. This is a qualitative lower bound; neither a numerical bound nor the minimizing quartet is determined here. Quantifying it would give an experimental calibration benchmark.

The upper bound does not pass from dimension six to its multiples. Already $12=4\times3$ supports four product MUBs. Five are possible using $I_{12}$ and four bases $D_j(F_6\otimes F_2)$, where each $D_j$ is diagonal with unit-modulus entries that depend jointly on both factor indices~\cite{wuttig2026mutuallyunbiasedbases}. These bases need not split into bases of the two factors. Restricting them to six dimensions preserves neither orthonormality nor unbiasedness in general. The order-six corner identities therefore impose no three-basis upper bound in dimension twelve.

For higher dimensions, we propose removing phase and permutation redundancies and seeking corner identities that express obstructions as sums of nonnegative terms, with certified covers of any remaining regions, including singular cases. Comparable reductions to the order-six reconstruction and maximum orbit would be needed before this becomes a practical method in dimensions ten and twelve.

Dimension six shows that the tensor product count can be exact, while dimension twelve shows that phases coupling the factors can overcome it. What property of a composite dimension determines which of these two possibilities occurs?

\section*{Acknowledgements}
M.C.W. acknowledges the hospitality of the Center for Computational Quantum Physics at the Flatiron Institute, a division of the Simons Foundation.

\paragraph*{AI assistance and author contributions.}
OpenAI's ChatGPT 6 Astra and Claude Fable 5.1 assisted in developing and checking candidate arguments. M.C.W. and J.T. directed the calculations, developed and reviewed the mathematical arguments, and prepared the manuscript. The authors are responsible for its content.

\bibliography{ref}

\clearpage
\onecolumngrid
\appendix
\section{Exact constellations}
\label{app:constellations}
A constellation $\{n_0,n_1,n_2,n_3\}_6$ consists of four groups of $n_j$ orthonormal vectors in $\C^6$, with squared overlaps $1/6$ between groups. The groups may be incomplete: these are not four full bases. Figure~\ref{fig:constellations} gives independent examples, not inputs to the exclusion proof. Numerical searches, including Douglas--Rachford iterations between Gram-matrix rank and overlap constraints, supply candidates; the analytical certificates in SM Sec.~IV establish these examples exactly. On an unmasked companion chart, $u_\pm$ are the unit-modulus row-product roots of the horizontal companion quadratic, and $\Delta u=|u_+-u_-|$~\cite{wuttig2026completeclassificationcomplexhadamard}.
\begin{figure}[h!]
 \centering
 \includegraphics[width=\textwidth]{Figure_3.jpg}
 \caption{\textbf{Three exact constellations of partial vector groups in dimension six.} \textbf{a,} $\{5,3,3,3\}_6$. \textbf{b,} $\{5,4,3,1\}_6$. \textbf{c,} $\{5,5,2,2\}_6$. The canonical basis is fixed in every panel; five basis rays determine the sixth. The compact diagrams show the remaining vector groups. Vectors within a group are orthogonal, and squared overlaps between groups are $1/6$. The pairs $B,C$ in \textbf{c} are generated by an antiunitary $K$ with $K^2=-I$. Each large torus retains $(\theta_1,\theta_2)$, and the smaller tori retain $(\theta_3,\theta_4)$ for selected dephased three-by-three corners. Matching point labels and shaded links identify the two coordinate pairs of each corner. Surface color gives the companion-root separation $\Delta u$ defined above, rather than the scalar $R$ of Fig.~\ref{fig:proof}; gray denotes the chart mask. Each large surface fixes the last two angles at its blue selected corner, while each small surface fixes its own first two angles. Other large-torus markers are projections. The chart mask alone does not establish existence: the full vectors and their relative phases are specified and certified exactly in SM Sec.~IV.}
 \label{fig:constellations}
\end{figure}
